\documentclass[10pt]{article}
\usepackage[margin=20mm]{geometry}
\usepackage{amsmath,amssymb,mathtools,fontspec,unicode-math,xurl,hyperref}
\IfFontExistsTF{DejaVu Serif}{\setmainfont{DejaVu Serif}}{\setmainfont{Latin Modern Roman}}
\setmathfont{latinmodern-math.otf}
\setlength{\parindent}{0pt}\setlength{\parskip}{6pt}\sloppy
\title{Verification of quartic representation arguments}
\author{AI-assisted technical review}
\date{}
\begin{document}
\maketitle
Abstract: The reviewed arXiv version contains an invalid lattice-point selection in the proof of Proposition 4.1. An explicit quartic satisfies its hypotheses and is nonnegative on every integer point, although the proof’s discriminant branch would force arbitrarily negative values. This is a Major defect in a core proof branch, not a counterexample to the principal nonrepresentation theorem. Two more limited errors in auxiliary statements are documented. These conclusions concern the exact 2023 preprint; correspondence with the later journal text has not been established.\par
\section*{Identity, version and inspected scope}
This report concerns Stanley Yao Xiao and Shuntaro Yamagishi, Quartic polynomials in two variables do not represent all non-negative integers, arXiv:2307.05712v2, dated 13 July 2023 [paper]. All page references are one-based PDF pages, agreeing with printed pages 1–25. The associated journal DOI is 10.4171/JEMS/1697 (https://doi.org/10.4171/JEMS/1697). The publisher records acceptance on 5 July 2025 and online publication on 6 August 2025 [publisher]; the journal PDF was not acquired. The findings must not be attributed to unchanged journal text.\par
The entire supplied PDF was read as extracted text. Detailed checking concentrated on Sections 2–3, the linear-factor case in Section 4, and the rational-curve case in Section 6.3. Section 5 and the remaining geometric case split were read for dependencies, without independently reproving the cited inhomogeneous approximation, Bernays, Siegel or quantitative Hilbert irreducibility results. No author program was run. Source-preparation searches and the present bounded public-record search located no separate supplement or paper-specific code; this availability limit is not a defect.\par
\section*{Framework and central dependency}
Write \(F=\sum_{i=1}^{4}F_i\in\mathbb Z[x,y]\), with \(F_i\) homogeneous of degree \(i\) and \(F(0,0)=0\), as in Section 1.1. Theorem 1.5 excludes \(F(\mathbb Z^2)=\mathbb Z_{\geq0}\) for rational quartics. Its proof reduces to Propositions 1.7–1.9; Proposition 1.9 in turn invokes Propositions 4.1, 4.2 and 5.5. Thus the alternative between sparse positive values and unbounded negative values in Proposition 4.1 is part of the core proof. A failed intermediate inference does not by itself refute this alternative or either main theorem.\par
\section*{Substantive findings}
Problem 1 (The unbounded-discriminant branch selects points outside the lattice).  Assessment: Major.\par
\subsection*{Source claim and location}
After (4.2), p. 13, the proof writes \[G(u,v)=g_2(u)v^2+g_1(u)v+g_0(u),\qquad
 D(u)=\frac{g_1(u)^2-4g_0(u)g_2(u)}{4g_2(u)}.\] Immediately after (4.3), it asserts that the quotient \(q\) and remainder \(h\) in \(g_1=2g_2q+h\), \(\deg h\leq1\), have integer coefficients. The first paragraph of p. 14 then takes \(v=-q(u)\) when \(D(u)\to+\infty\), obtains \(G(u,v)=-D(u)+O(1)\), and concludes (4.6), unboundedness below on \(\mathbb Z^2\). Division over \(\mathbb Q[u]\) does not justify that integer-point selection.\par
\subsection*{An admissible quartic}
Equation reference: eq:F. Set \[
 Q(u,v)=u^2+(2u+1)v,\qquad F(u,v)=16Q(u,v)^2-u(u+1).\] This integer quartic has \(F(0,0)=0\) and homogeneous parts \[\begin{aligned}
 F_4&=16u^2(u+2v)^2,&F_3&=32uv(u+2v),\\
 F_2&=16v^2-u^2,&F_1&=-u.
 \end{aligned}\] In Proposition 4.1 choose \(H=u\), \(G_4=16(u+2v)^2\) and \(G_3=32v(u+2v)\). The form \(H\) is primitive and linear; \(G_4,G_3\) are integral quadratic forms, neither divisible by \(H\). The two linear factors of \(G_3\) are distinct, so the restriction excluding a square linear divisor not proportional to \(H\) is satisfied. Also \(F_4\geq0\).\par
The example is linked: \(\widetilde F_4=u^2(u+2v)^2\), \(\gcd(\widetilde F_4,F_3)=u(u+2v)\), and its gcd with \(F_2\) is 1. On both zero directions of \(F_4\), \(F_2\) is strictly positive. Both cubic multiplicities are 1 and both quartic multiplicities are 2. Thus neither an earlier negative quadratic direction nor unequal multiplicity excludes this example.\par
Equation reference: eq:division. The identity change of variables gives \[\begin{aligned}
 g_2&=16(2u+1)^2,&g_1&=32u^2(2u+1),\\
 g_0&=16u^4-u(u+1),&D(u)&=u(u+1).
 \end{aligned}\] Here \(g_2(u)>0\) for every integer \(u\), and \(D(u)\to+\infty\). The actual quotient and remainder are \[
 q(u)=\frac u2-\frac14,\qquad h(u)=16u+8.\] Uniqueness of polynomial division rules out another integral quotient with remainder of degree at most 1. Moreover \(q(u)=(2u-1)/4\notin\mathbb Z\) for every integer \(u\). The proposed evaluation does give \[F(u,-q(u))=1-u(u+1),\] but only at points outside \(\mathbb Z^2\).\par
\subsection*{Universal nonnegativity on integer points}
If \(u=2k\geq0\), then \[Q(u,v)=(4k+1)(v+k)-k,\qquad |Q(u,v)|\geq k=u/2.\] If \(u=2k+1\geq1\), then \[Q(u,v)=(4k+3)(v+k)+(k+1),\qquad |Q(u,v)|\geq k+1.\] These bounds are the least absolute residues modulo the displayed positive odd moduli. Consequently, for \(u\geq1\), \[F(u,v)\geq4u^2-u(u+1)=3u^2-u>0.\] For \(u=0\), \(F(0,v)=16v^2\geq0\). Finally, \[Q(-u-1,1-v)=Q(u,v),\qquad F(-u-1,1-v)=F(u,v).\] This involution covers all negative integers \(u\). Hence \(F\geq0\) on \(\mathbb Z^2\) and its infimum is 0, contrary to the asserted consequence of the \(D(u)\to+\infty\) branch.\par
\subsection*{Impact and repair boundary}
An implicit nonnegativity assumption cannot exclude the example: it satisfies that stronger condition. Replacing integer by rational in the division statement does not repair the subsequent lattice-point selection. Rounding does not retain an \(O(1)\) error when \(g_2(u)\) grows quadratically. For example, at \(u=20\), the rational point gives \(-419\), whereas the minimum over integer \(v\) is \(1180\), at \(v=-10\).\par
This demonstrated invalid inference lies in the core linear-factor branch and supports Major. It does not prove Proposition 4.1 or either main theorem false. Indeed [eq:F] (\#eq:F) takes only even values and is not a counterexample to the nonrepresentation theorem. A replacement argument must control the distance from the rational center to the lattice relative to \(D(u)\). No general repair of that branch is verified here.\par
Problem 2 (An angular lower bound omits unshared quartic zero directions).  Assessment: Minor; literal auxiliary-statement correction.\par
Section 3, p. 6, defines \(\Xi\) using only common real zero directions of \(\widetilde F_4\) and \(F_3\). Lemma 3.1, p. 7, nevertheless claims \(F(B\cos\theta,B\sin\theta)\geq C_1B^3\) outside the shrinking neighborhoods of \(\Xi\).\par
Take \(G(x,y)=x^2y^2+x^3+y^2\). Then \(G_4\geq0\), \(\gcd(\widetilde G_4,G_3)=x^2\neq1\), and \(\gcd(\widetilde G_4,G_3,G_2)=1\); the polynomial is linked. Its set \(\Xi\) consists of the two directions \(x=0\), omitting \(y=0\). For \(n\geq1\), let \[(x_n,y_n)=(-n^2,n),\qquad B_n=\sqrt{n^4+n^2}.\] These directions remain a positive angular distance from \(\Xi\), hence lie outside its \(O(B_n^{-1/2})\) neighborhoods for large \(n\), but \[G(x_n,y_n)=n^2>0,\qquad
 \frac{G(x_n,y_n)}{B_n^3}
 =\frac1{n(n^2+1)^{3/2}}\longrightarrow0.\] Thus even restricting to positive values does not suffice.\par
A sufficient local correction requires every real zero direction of \(F_4\) to belong to \(\Xi\). Near a zero of multiplicity \(2r\), outside a \(cB^{-1/(2r)}\) neighborhood, the quartic term is bounded below by a constant times \(c^{2r}B^3\); choosing \(c\) large dominates the cubic term. Compactness treats the remaining directions. A globally bounded-below quartic cannot have an unshared zero direction, by applying the approximation and sign argument of Lemma 2.5 at that direction.\par
Here \(G(-n,0)=-n^3\), so the example already satisfies the negative-value alternative. It establishes a missing hypothesis in the auxiliary lemma, not damage to the intended bounded-below target. That narrower impact, rather than the ease of repair, motivates Minor. The invocation in Proposition 3.2 needs the stated preliminary reduction.\par
Problem 3 (A constant transformed denominator need not have constant numerator).  Assessment: Minor; omitted rational-polynomial case.\par
Section 6.3, p. 24, asserts \(\deg\widetilde Q_2=0\) if and only if \(\deg R_2=\deg Q_2=0\). Take \(F(x,y)=x(x^3-y+1)\), \(K=x^3-y+1\) and \(d_2=1\). Its leading form is positive semidefinite and \(\deg\gcd(\widetilde F_4,F_3,F_2,F_1)=1\), as in Lemma 6.8. The curve \(K-d_2=x^3-y=0\) is geometrically irreducible of genus zero. Use \[R_1=t,\quad Q_1=1,\quad R_2=t^3,\quad Q_2=1.\] In (6.2), \(r=v=1\), \(s=u=0\), and \(m=3\). The definitions give \(\widetilde R_2=-x^3\), \(\widetilde Q_2=-1\). Thus the denominator is constant although \(R_2\) has degree 3.\par
This does not refute Lemma 6.8: \(F(x,x^3)=x\) is already unbounded below. The omitted case is \(y=P(x)\) for a nonconstant rational polynomial. If \(P(x_0)\in\mathbb Z\), clearing coefficient denominators supplies an arithmetic progression \(x=x_0+Mk\) with \(P(x)\in\mathbb Z\), unbounded in both signs. This addresses the polynomial-graph subcase, not all the geometric reductions.\par
\section*{Overall assessment}
The highest supported impact is Major, for Problem 1. No counterexample to the main theorem or complete corrected proof is claimed. Problems 2–3 are separately graded because their examples already satisfy the negative-value alternative. The journal text may differ; no comparison establishes whether these passages were corrected.\par
Item | Assessment | Repair verification\par
Problem 1 | Major | General repair not verified\par
Problem 2 | Minor | Sufficient local condition proved\par
Problem 3 | Minor | Polynomial-graph subcase addressed\par
Scientific decision | Pending | Human review not supplied\par
Author adoption | Unknown | No author contact made\par
\section*{Evidence, calculations and limitations}
Source PDF SHA256: 03eef10ef03e0229186459a6b3e12506472b77cd7d1172af7aad8a5b088f7538. Independent standard-library Python checked coefficient identities, the discriminant identity and polynomial division with exact fractions. A diagnostic over \(80{,}601\) integer pairs found minimum 0. The universal proof is the parity-and-involution argument, not that grid. Code, actual command, exit code and output are retained under `verification/0321/checks/`.\par
In the same-conversation continuation, the nine source evidence pages were actually displayed and inspected through the authorized native image route. This revision updates that operational status only; the scientific findings are unchanged. Exact compilation, final-page observations and image hashes are recorded separately in the current provenance. The controller supplied the UI observation; this worker did not observe it. Runtime self-description and UI attribution remain separate. No human review, independent acceptance or journal endorsement is claimed.\par
9 S. Y. Xiao and S. Yamagishi, Quartic polynomials in two variables do not represent all non-negative integers, arXiv:2307.05712v2, 13 July 2023, https://arxiv.org/abs/2307.05712v2. EMS Press, bibliographic record, DOI 10.4171/JEMS/1697, accessed 27 October 2026, https://ems.press/journals/jems/articles/14299022. Metadata only; journal PDF not reviewed.\par
\end{document}
